\chapter*{General Introduction}
\addcontentsline{toc}{chapter}{General Introduction}
\markboth{General Introduction}{General Introduction}

\section*{Context of the Internship}
\addcontentsline{toc}{section}{Context of the Internship}

This thesis documents a six-month Master's internship project, carried out as part of the Master's Program in Data Sciences at Polytech Intl, addressing a concrete, everyday operational problem for a real client organization, the African Development Bank Group: the Bank's Delegation of Authority Matrix (DAM) -- the internal document defining who may approve, check, initiate, or must be informed about each of the Bank's operational activities -- is authoritative but practically unreadable at speed. It is distributed as a 79-page PDF of dense, rotated-header tables, footnoted role names, and a bespoke authority-code notation. Anyone who needs an answer from it -- a new staff member, a manager double-checking a delegation, an auditor tracing accountability for a decision -- must manually locate the right table, read across a 90-degree-rotated header row, and cross-reference footnotes scattered across the page. There is no search function smarter than literal string matching in a PDF reader, and no way to ask the document a question in plain language.

The internship's mandate was to determine whether this problem could be solved responsibly with an AI system: not a generic chatbot bolted onto the PDF, but a system whose answers could be trusted for a compliance-sensitive purpose, where a wrong answer about who has authority to approve a transaction is a real institutional risk, not a minor inconvenience. That framing -- trustworthiness first, conversational convenience second -- is the thread running through every technical decision documented in this thesis.

\section*{Report Structure and How to Read It}
\addcontentsline{toc}{section}{Report Structure and How to Read It}

This thesis is organized to serve two audiences at once: a reader who wants the engineering narrative (what was built, what went wrong, how it was fixed, what remains open) and a reader who wants to verify specific claims against primary evidence (exact test counts, exact bug descriptions, exact source-code behavior). To serve both without forcing either to wade through the other's material, the thesis is split into a main body and a substantial appendix section.

Chapter~1 presents the host institution, Polytech Intl. Chapter~2 presents the client organization, the African Development Bank Group, and the working relationship between the two. Chapter~3 provides the theoretical and methodological background needed to follow the technical decisions in later chapters -- CRISP-DM itself, knowledge graphs, retrieval-augmented generation, and the specific extraction and search techniques used. Chapters~4 through~10 are the CRISP-DM phases proper: Business Understanding, Requirements Specification, Data Understanding, Data Preparation, Modeling, Evaluation, and Deployment. Chapter~11 is a practical user manual for anyone who wants to run or extend the system themselves. Chapters~12 and~13 close the main body with project-management reflection and an honest accounting of limitations, followed by a General Conclusion and a Web References section.

The appendices are deliberately extensive and are not optional reading for anyone verifying this thesis's claims rather than just following its narrative. Annex~F reproduces the project's complete, dated decision log in full -- every bug found, every fix attempted (including the ones that failed and were reverted), every scope decision, in the order it actually happened. It is the single primary source this entire thesis is built from, and it is included in full rather than only excerpted, in the same spirit that a scientific report includes raw data rather than only summary statistics. Annex~G reproduces the source code of the system's most technically significant functions directly, so that a claim like ``the grounding check normalizes whitespace and case before comparing'' can be checked against the actual sixty lines of Python that implement it, not just a paraphrase of them.

\section*{A Note on Authorship and Process}
\addcontentsline{toc}{section}{A Note on Authorship and Process}

This project was carried out by one graduate intern working in close, continuous collaboration with an AI coding assistant acting in a senior-engineer advisory role -- reviewing proposed approaches, flagging risks before code was written, and, critically, insisting that every claimed fix be checked against real data before being trusted. That collaboration is itself part of the story this thesis tells: the discipline of ``verify against a real screenshot or real characters before touching code, and honestly document a fix attempt that failed rather than silently discarding it'' is visible throughout the decision log in Annex~F, and this thesis treats that discipline as one of the project's genuine deliverables, not merely a process footnote.
